%% kapa: exam=1 round=24 year=2013 period=1 form=A subject=civil_law subject_ko=민법 seq=24 number=24
%% source: Q-Net 감정평가사 기출문제 (cst00309) artlSeq=5073802 file=extracted/1008/5073802/2013년 제24회 감정평가사 1차시험 기출문제/1교시(A형).hwp
%% conversion: kordoc 4.24.51 -> tools/gen1.py (auto); review=proofread-sonnet; flags=none
\begin{question}{24}
청구권 보전의 가등기에 관한 설명으로 옳지 않은 것은? (다툼이 있으면 판례에 의함)
\begin{choices}
  \item 본등기의 순위는 가등기의 순위에 따르나 물권변동의 시기가 가등기를 경료한 때로 소급하지 않는다.
  \item 소유권이전등기청구권 보전을 위한 가등기가 있다 하여 소유권이전등기를 청구할 어떤 법률관계가 있다고 추정되지 않는다.
  \item 소유권이전등기청구권의 보전을 위한 가등기에 기한 본등기가 경료되면 그 사이에 행해진 중간처분의 등기는 가등기에 의하여 순위가 보전된 권리와 양립할 수 없는 범위 내에서 무효로 된다.
  \item 매수인 명의의 가등기가 이루어진 후에 매도인이 제3자 앞으로 소유권이전등기를 경료해 준 경우 매수인은 매도인을 상대로 본등기를 청구하여야 한다.
  \item 부동산 강제경매절차에서 부동산이 매각된 경우 특단의 사정이 없으면 소유권이전등기청구권의 보전을 위한 가등기는 말소된다.
\end{choices}
\end{question}
\answer{5}
